\makeglossaries
\newglossaryentry{Laue's theory}
{
    name=Laue's theory,
    description={Laue's Law is a fundamental principle in X-ray diffraction used to describe the scattering of X-rays by a crystal lattice. It states the conditions that must be met for constructive interference to occur when a beam of X-rays interacts with the regularly spaced atoms in a crystal.In terms of the reciprocal lattice, the condition for diffraction is that the difference between the incident wave vector ($\mathbf{k}$) and the scattered wave vector ($\mathbf{k'}$) must equal a reciprocal lattice vector ($\mathbf{G}$).}
}

\newglossaryentry{Laue's equation}{
    name= Laue's equation,
    description= {The three Laue Equations are the scalar expressions that must be simultaneously satisfied for constructive interference to occur when X-rays are diffracted by the periodic structure of a crystal lattice.They represent the condition that the path difference between waves scattered from adjacent atoms along each of the three crystal axes ($\mathbf{a}$, $\mathbf{b}$, $\mathbf{c}$) must be an integral multiple of the X-ray wavelength ($\lambda$).The equations are:$$a(\cos\alpha - \cos\alpha_0) = h\lambda \\
b(\cos\beta - \cos\beta_0) = k\lambda \\
c(\cos\gamma - \cos\gamma_0) = l\lambda$$Where:$a, b, c$ are the primitive translation lengths (lattice constants) along the crystal axes.$\cos\alpha_0$, $\cos\beta_0$, $\cos\gamma_0$ are the direction cosines of the incident X-ray beam with respect to the crystal axes.$\cos\alpha$, $\cos\beta$, $\cos\gamma$ are the direction cosines of the diffracted X-ray beam with respect to the crystal axes.$h, k, l$ are integers (now known to be equivalent to the Miller indices of the diffracting plane).$\lambda$ is the X-ray wavelength.These three scalar equations are mathematically equivalent to the single vector form of Laue's Law, which states that the change in the wave vector must equal a reciprocal lattice vector $\mathbf{G}$:$$\mathbf{k'} - \mathbf{k} = \mathbf{G}$$}
}

\newglossaryentry{Miller Indices}{
    name= Miller indices,
    description={Miller Indices are a system of notation in crystallography that uses three integers to uniquely identify the orientation of a family of parallel planes within a crystal lattice.They are derived from the reciprocals of the fractional intercepts that a plane makes with the three crystallographic axes ($\mathbf{a}$, $\mathbf{b}$, $\mathbf{c}$).A set of Miller indices $(hkl)$ for a plane is fundamentally equivalent to identifying a vector in the reciprocal lattice, $\mathbf{G}_{hkl}$, which is perpendicular to that plane.The plane's equation, in terms of its fractional intercepts $(a/h, b/k, c/l)$ along the crystal axes, satisfies:$$\frac{x}{a/h} + \frac{y}{b/k} + \frac{z}{c/l} = n \quad \text{or, more generally, } \quad h\frac{x}{a} + k\frac{y}{b} + l\frac{z}{c} = n$$where:$(hkl)$ are the Miller Indices (always integers, often reduced to the smallest set).$n$ is an integer (the $n$-th plane from the origin).A bar over an index (e.g., $\bar{h}$) denotes a negative intercept.Round brackets $\mathbf{(hkl)}$ denote a single plane or a family of parallel planes.Curly brackets $\mathbf{\{hkl\}}$ denote a family of equivalent planes (by crystal symmetry).They are essential for interpreting X-ray diffraction patterns, as the diffraction condition (Laue's Law) directly relates the change in the wave vector to a reciprocal lattice vector $\mathbf{G}_{hkl}$, which is indexed by $h, k, l$.}
}